% Блок-схема генетического алгоритма подбора параметров политики охлаждения (GA/ga_core.py).
% Компиляция: pdflatex ga_block_diagram.tex
%
% Редактирование: TeXstudio, VS Code + LaTeX Workshop, Overleaf.

\documentclass[12pt,a4paper,landscape]{article}
\usepackage[utf8]{inputenc}
\usepackage[T2A]{fontenc}
\usepackage[russian]{babel}
\usepackage{amsmath,amssymb}
\usepackage{tikz}
\usetikzlibrary{positioning,arrows.meta,calc,shapes.geometric}

\usepackage[margin=1.2cm]{geometry}

\title{Блок-схема генетического алгоритма настройки политики охлаждения}
\author{}
\date{}

\begin{document}
\maketitle

\noindent
\textbf{Назначение:} минимизация приспособленности
$F = E_{\mathrm{cool}} + w \cdot P_{\mathrm{temp}}$,
где $E_{\mathrm{cool}}$~--- интеграл мощности охлаждения по эпизоду симуляции,
$P_{\mathrm{temp}}$~--- штраф за превышение средней температуры чипов над лимитом
(см.\ \texttt{ga\_core.py}: \texttt{evaluate\_chromosome}).

\vspace{0.5em}

% --- Рис. 1: run_ga ---
\begin{center}
\begin{tikzpicture}[
  font=\footnotesize,
  >=Stealth,
  blk/.style={
    rectangle,
    draw=black!70,
    rounded corners=3pt,
    thick,
    align=center,
    inner sep=6pt,
    minimum width=3.4cm,
  },
  dec/.style={
    diamond,
    draw=black!70,
    thick,
    aspect=2.2,
    inner sep=2pt,
    align=center,
  },
  smallblk/.style={
    rectangle,
    draw=black!55,
    rounded corners=2pt,
    align=left,
    inner sep=5pt,
    font=\scriptsize,
  },
]

  \node[blk, fill=green!12] (start) {\textbf{Начало}\\[0.2em]задать ГПСЧ};

  \node[blk, fill=blue!8, below=0.55cm of start] (init) {%
    \textbf{Инициализация популяции}\\[0.2em]
    $N_{\mathrm{pop}}$ случайных хромосом\\[0.1em]
    {\scriptsize гены: $T^{\mathrm{tgt}}_{\mathrm{chip}}$, $k_p$, $s^{\mathrm{base}}_{\mathrm{fan}}$,
    $k_{\mathrm{fan}}$, $\Delta sp_{\max}$}};

  \node[blk, fill=yellow!10, below=0.5cm of init] (loop) {%
    \textbf{Цикл по поколениям}\\[0.15em]
    $g = 1, 2, \ldots, G$};

  \node[blk, fill=orange!14, below=0.5cm of loop, minimum width=4.2cm] (fitall) {%
    \textbf{Оценка популяции}\\[0.15em]
    для каждой хромосомы $i$:\\[0.1em]
    $f_i \leftarrow F(g_i)$ \quad (прогон API двойника)};

  \node[blk, fill=purple!10, below=0.5cm of fitall] (best) {%
    \textbf{Лучший в поколении}; обновить\\[0.1em]
    глобально лучшую хромосому};

  \node[blk, fill=cyan!10, below=0.5cm of best, minimum width=4.8cm] (newpop) {%
    \textbf{Новое поколение}\\[0.15em]
    элита: копия лучшего индивида;\\[0.1em]
    дополнить до $N_{\mathrm{pop}}$:\\[0.15em]
    турнир $(k{=}3)$ $\rightarrow$ кроссовер $\rightarrow$ мутация};

  \node[blk, fill=gray!18, below=0.55cm of newpop] (end) {%
    \textbf{Конец}\\[0.15em]
    вернуть лучшую хромосому за все $G$ поколений};

  \draw[thick,->] (start) -- (init);
  \draw[thick,->] (init) -- (loop);
  \draw[thick,->] (loop) -- (fitall);
  \draw[thick,->] (fitall) -- (best);
  \draw[thick,->] (best) -- (newpop);
  \draw[thick,->] (newpop) -- (end);

  % Пояснение справа
  \node[smallblk, fill=white, right=1.1cm of fitall, anchor=north west, yshift=0.3cm] (note1) {%
    \textbf{Детализация} $F(g_i)$ --- рис.~2.\\[0.2em]
    Кроссовер: взвешенное смешивание\\[0.1em]
    родителей, $\alpha \in [0{,}2,\,0{,}8]$.\\[0.2em]
    Мутация: гауссов шум по гену\\[0.1em]
    с вероятностью $0{,}15$ (после клипа в границы).
  };

\end{tikzpicture}
\end{center}

\vspace{0.4em}

% --- Рис. 2: evaluate_chromosome ---
\begin{center}
\begin{tikzpicture}[
  font=\footnotesize,
  >=Stealth,
  blk/.style={
    rectangle,
    draw=black!70,
    rounded corners=3pt,
    thick,
    align=center,
    inner sep=6pt,
    minimum width=3.6cm,
  },
  dec/.style={
    diamond,
    draw=black!70,
    thick,
    aspect=2.4,
    inner sep=1pt,
    align=center,
    font=\scriptsize,
  },
]

  \node[blk, fill=green!12] (e0) {\textbf{Вход}: хромосома $g$};

  \node[blk, fill=blue!8, below=0.45cm of e0] (e1) {%
    \textbf{Сценарий на двойнике}\\[0.15em]
    сброс, режим охлаждения, уставка/вент.,\\[0.1em]
    нагрузка, $T_{\mathrm{out}}$};

  \node[blk, fill=yellow!10, below=0.45cm of e1] (e2) {%
    $s \leftarrow 1$; накопители:\\[0.1em]
    $E_{\mathrm{cool}} \leftarrow 0$, $P_{\mathrm{temp}} \leftarrow 0$};

  \node[blk, fill=orange!12, below=0.45cm of e2, minimum width=4cm] (e3) {%
    \textbf{Шаг симуляции} ($s$)\\
    \texttt{/simulation/step}; телеметрия};

  \node[blk, fill=orange!8, below=0.4cm of e3] (e4) {%
    учёт $E_{\mathrm{cool}} += P_{\mathrm{inf}}\,\Delta t$;\\[0.1em]
    штраф по $\overline{T_{\mathrm{chip}}}$};

  \node[dec, fill=magenta!6, below=0.45cm of e4] (e5) {шаг кратен\\[0.1em]\texttt{control\_interval}?};

  \node[blk, fill=cyan!12, below left=0.55cm and 0.05cm of e5] (e6) {%
    P-закон по ошибке\\[0.1em]
    $e = \overline{T_{\mathrm{chip}}} - T^{\mathrm{tgt}}_{\mathrm{chip}}$:\\[0.1em]
    $\Delta T_{\mathrm{set}}$, новая уставка;\\[0.1em]
    $s_{\mathrm{fan}}$ из $g$};

  \node[blk, fill=gray!15, below right=0.55cm and 0.05cm of e5] (e7) {без изменения\\[0.1em]управления};

  \coordinate (join) at ($(e6.south)!0.5!(e7.south)$);
  \node[dec, below=0.65cm of join] (e8) {%
    $s < n_{\mathrm{steps}}$?\\[0.15em]
    {\scriptsize (ещё шаг эпизода?)}};

  \node[blk, fill=purple!12, below=0.65cm of e8] (e9) {%
    \textbf{Выход}: $F = E_{\mathrm{cool}} + w\,P_{\mathrm{temp}}$};

  \draw[thick,->] (e0) -- (e1);
  \draw[thick,->] (e1) -- (e2);
  \draw[thick,->] (e2) -- (e3);
  \draw[thick,->] (e3) -- (e4);
  \draw[thick,->] (e4) -- (e5);
  \draw[thick,->] (e5) -- node[left, font=\scriptsize, pos=0.4] {да} (e6);
  \draw[thick,->] (e5) -- node[right, font=\scriptsize, pos=0.4] {нет} (e7);
  \draw[thick,->] (e6) |- (join);
  \draw[thick,->] (e7) |- (join);
  \draw[thick,->] (join) -- (e8);
  % да: ещё есть шаги — снова шаг симуляции (следующий s)
  \draw[thick,->] (e8.east) -- node[above, font=\scriptsize, near start] {да} ++(1.85,0)
    |- node[right, font=\scriptsize, pos=0.22] {$s \leftarrow s{+}1$} (e3.east);
  \draw[thick,->] (e8.south) -- node[right, font=\scriptsize, pos=0.42] {нет} (e9);

\end{tikzpicture}
\end{center}

\vspace{-0.8em}
\begin{center}
\footnotesize
\textit{Рис.~2.} После обработки шага $s$ проверяется, остались ли шаги эпизода;
при \textbf{да} выполняется $s \leftarrow s{+}1$ и снова вызывается шаг симуляции.
\end{center}

\vspace{0.3em}

\noindent\textbf{Соответствие коду.}
\begin{align*}
  &\textbf{Гены} \quad g = (T^{\mathrm{tgt}}_{\mathrm{chip}},\, k_p,\, s^{\mathrm{base}}_{\mathrm{fan}},\, k_{\mathrm{fan}},\, \Delta sp_{\max}).
  \\[0.35em]
  &\textbf{Приспособленность.} \quad
  F = \sum_{s=1}^{n_{\mathrm{steps}}} P_{\mathrm{inf}}(s)\,\Delta t
  + w \sum_{s=1}^{n_{\mathrm{steps}}} \max\bigl(0,\, \overline{T_{\mathrm{chip}}}(s) - T_{\mathrm{lim}}\bigr)\,\Delta t.
  \\[0.35em]
  &\textbf{Управление} (на шагах кратных \texttt{control\_interval}):
  \quad
  \Delta T_{\mathrm{set}} = -\mathrm{clip}(k_p\, e,\, \pm \Delta sp_{\max}), \quad
  s_{\mathrm{fan}} = \mathrm{clip}\bigl(s^{\mathrm{base}}_{\mathrm{fan}} + k_{\mathrm{fan}}\max(0,e),\, 0,\,100\bigr).
\end{align*}

\end{document}
